
\paragraph{\textbf{Corollary 2.14.}}
\SGARefTarget{sga3:xiv:2.14}
\emph{When condition \((\mathrm C_3)\) of \SGARefLink{sga3:xiv:2.9}{2.9} holds, the Cartan
subalgebras of \(\mathfrak g\) are strictly nilpotent
\textup{(\SGARefLink{sga3:xiv:2.2}{2.2})}.}

The proof is immediate.

\paragraph{\textbf{Remark 2.15.}}
\SGARefTarget{sga3:xiv:2.15}
\noindent
\textup{a)}
One can prove a converse to \SGARefLink{sga3:xiv:2.14}{2.14}: condition \((\mathrm C_3)\) is
equivalent to the assertion that, for every \(S'\) over \(S\),
every quasi-regular section of \(\mathfrak g_{S'}\) is regular and
every Cartan subalgebra of \(\mathfrak g_{S'}\) is strictly
nilpotent; equivalently, every quasi-regular section of
\(\mathfrak g_{S'}\) is contained in a strictly nilpotent Cartan
subalgebra of \(\mathfrak g_{S'}\).

\smallskip
\noindent
\textup{b)}
Unlike the other conditions \((\mathrm C_0)\) through
\((\mathrm C_2)\), condition \((\mathrm C_3)\) is of
``infinitesimal nature.'' More precisely, when \(S\) is locally
noetherian, \(\mathfrak g\) satisfies \((\mathrm C_3)\) if and only
if, for every base change \(S'\to S\) with \(S'\) local Artinian,
\(\mathfrak g_{S'}\) satisfies the same condition. If desired, one
may take \(S'\) to be the spectrum of an Artinian quotient of a
local ring of \(S\). Similarly, when \((\mathrm C_0)\) holds, the
condition for a section of \(\mathfrak g\) to be regular is of
infinitesimal nature.

\smallskip
\noindent
\textup{c)}
When \(\mathfrak g\) is the Lie algebra of a smooth group
prescheme of finite presentation over \(S\), we shall see that the
conditions
\[
  (\mathrm C_0),\quad(\mathrm C_1),\quad
  (\mathrm C'_1),\quad(\mathrm C_2)
\]
on \(\mathfrak g\) are equivalent \textup{(\SGARefLink{sga3:xiv:5.2}{5.2}~a)}. I do not know
what happens in general, except that even when \(S\) is local
Artinian, \((\mathrm C_0)\) does not imply \((\mathrm C_1)\).
However, even when \(\mathfrak g\) comes from a group \(G\) and
\(S\) is local Artinian, it is not true in general that
\((\mathrm C_2)\) implies \((\mathrm C_3)\), since
\(\mathfrak g\) may be nilpotent without being strictly
nilpotent.

An example is obtained from a smooth affine group scheme \(G\)
over the spectrum \(S\) of a discrete valuation ring, such that
the Lie algebra of the generic fiber is not nilpotent---for
example, the generic fiber may be an adjoint semisimple
group---whereas that of the special fiber is nilpotent---for
example, the special fiber may be a vector group. Then, for
sufficiently large \(n\), the Lie algebra of
\[
  G_n=G\times_S S_n
\]
is not strictly nilpotent, although it is nilpotent.
